\documentclass[10pt,a4paper]{amsart}
\usepackage[margin=19mm]{geometry}
\usepackage{amsmath,amssymb,mathtools}
\usepackage[scr=rsfs,cal=euler]{mathalfa}
\usepackage{xfrac}
\usepackage[unicode,hidelinks,bookmarks=false]{hyperref}
\newcommand{\ie}{i.e.\ }
\newcommand{\cf}{cf.\ }
\newcommand{\resp}{respectively\ }
\newcommand{\Subsection}{\subsection}
\providecommand{\nobreakdash}{\nobreak\hbox{-}}
\setlength{\emergencystretch}{2em}
\multlinegap0pt
\usepackage{amssymb}
\usepackage{bm}
\usepackage{braket}
\usepackage{xcolor}
\usepackage{tikz}
\newtheorem{lemma}{Lemma}
\usepackage{cleveref}
\crefname{lemma}{Lemma}{Lemmas}
\newcommand{\beq}{\begin{equation}}
\newcommand{\eeq}{\end{equation}}
\title{Proof of the Error Scaling for Universally Robust Dynamical Decoupling Sequences}
\author{Domenico D'Alessandro, Phattharaporn Singkanipa, Daniel Lidar}
\date{Source: arXiv:2604.25807v1 (CC BY 4.0)}
\begin{document}
\maketitle

\noindent\emph{Source and licence.} Proof of the Error Scaling for Universally Robust Dynamical Decoupling Sequences. Version arXiv:2604.25807v1 (CC BY 4.0), distributed by arXiv under the Creative Commons Attribution 4.0 International licence, http://creativecommons.org/licenses/by/4.0/, as recorded in the arXiv licence metadata for this exact version and shown on the abstract page https://arxiv.org/abs/2604.25807v1. Full licence notice: \texttt{licenses/arxiv-2604.25807-CC-BY-4.0.txt}.\par\smallskip

\noindent\textit{Editorial scope.} Counted unit: the article's lemma lem:odd\_reflection (source0.tex lines 933--940). The article's complete original proof of it is delivered with it (source0.tex lines 941--985, one \texttt{proof} environment of 45 lines). No mathematics is written, repaired or re-derived: the statement and the proof are the article's own lines, in the article's own notation, and no retained line is substituted or re-flowed.  No further source-local context is needed: every cross-reference of every retained line resolves inside the two delivered spans.  Nothing was omitted from any retained span, so the package carries no omission record.  No retained line contains a citation, so the wrapper carries no bibliography and no bibliography entry.  No retained line contains a figure environment, an image-inclusion command or drawing code, so no image is delivered and the package declares no asset dependency.  Editorial formatting only, in addition: an AMS \texttt{amsart} preamble that restates the article's own declarations and macros used by the retained lines, namely the environments \texttt{lemma} on separate counters, together with the standard packages those lines need.  The retained environments print in this document as Lemma 1 in order of appearance, whereas the article prints the same items with the numbering of its own full document.  The complete native source of the article as distributed in the arXiv e-print for this exact version is bundled unchanged as \texttt{sources/199330/source0.tex}.
\begin{lemma}\label{lem:odd_reflection}
For every $1 \leq k < \frac{n}{2}$ and every even $s$ with $0\leq |s|
 \leq 2k$ 
\beq
A_s^{(k)}=\omega^s (A_s^{(k)})^{*}
\label{eq1}
\eeq
\end{lemma}

\begin{proof}
For a given $k$, on the set of increasing sequences
\beq\label{Dset}
\mathcal D := \Bigl\{  r=(r_1,\dots,r_{2k})\in\mathbb Z^{2k}\ :\ 1\le r_1<\cdots<r_{2k}\le n \Bigr\}, 
\eeq
we define a map   $\tau   :   \mathcal D \rightarrow  \mathcal D$ componentwise as 
\beq
\tau(r)_j  :=  n+1-r_{2k+1-j}\quad (1\le j\le 2k).
\eeq
Direct verification shows that $\tau$ does indeed preserve the ordering in a sequence $r$ and it is a bijection, in fact an involution ($\tau^2$ is the identity map), since 
\beq
\tau(\tau(r))_j
=n+1-\bigl(n+1-r_{2k+1-(2k+1-j)}\bigr)=r_j. 
\eeq
We  now verify how the signature and weighted signature of $r$ and $\tau(r)$ are related. Using the definition of $\tau$, we have 
\beq 
\begin{aligned}
(-1)^j (-1)^{\tau(r)_j}&=(-1)^j(-1)^{n+1-r_{2k-j+1}}\\
=(-1)^{j+1}(-1)^{r_{2k-j+1}}&
=(-1)^{2k-j+1}(-1)^{r_{2k-j+1}}, 
\end{aligned}
\eeq
since $2k+1$ is odd. Summing over $j=1,...,2k$ and using a change of summation index $l=2k-j+1$, we have that $S(\tau(r))=S(r)$. 

The weighted signature $W$ transforms as follows:

\beq
\begin{aligned}
   W(\tau(r))
&= \sum_{j=1}^{2k}    (-1)^j (-1)^{\tau(r)_j} \tau(r)_j\\
 &=\sum_{j=1}^{2k}    (-1)^j (-1)^{\tau(r)_j}(n+1-r_{2k-j+1})\\
 &=(n+1)S(r)-\sum_{j=1}^{2k}(-1)^j (-1)^{n+1-r_{2k-j+1}}r_{2k-j+1}\\
 &=(n+1)S(r)-\sum_{l=1}^{2k}(-1)^l (-1)^{r_{l}}r_{l}\\
 &=(n+1)S(r)-W(r)
\end{aligned}
\eeq
where we again used the change of index $l=2k-j+1$. Therefore, using $\omega^n=1$, we have 
\beq
\omega^{W(\tau(r))}
= \omega^{(n+1)S(r)}\omega^{-W(r)}
= \omega^{S(r)}\big(\omega^{W(r)}\big)^*.
\eeq
Summing over all tuples with $S(r)=s$ and changing variables $r\mapsto \tau(r)$, we obtain \cref{eq1}. 

\end{proof}

\end{document}
